\section{Physical Interfaces}

\subsection{Timing Interface (DTS)}
\label{sec:physical-timing}

\paragraph{Physical layer.}
% Timing and synchronization are distributed over single-mode OS2 fiber using bidirectional 1000BASE-BX-20 SFPs:
% \begin{itemize}[nosep]
%   \item PDS timing endpoints: TX \SI{1550}{nm}, RX \SI{1490}{nm}, as specified by EDMS 2580777; the timing-master optics are complementary.
%   \item Connectors: LC on both ends; one LC fiber per endpoint (DAPHNE FEB and LCM).
%   \item Fan-out: last stage on the DAQ mezzanine or mini-racks; PDS owns fibers from the last fan-out to endpoints.
% \end{itemize}

% \paragraph{Reference and quality.}
% \begin{itemize}[nosep]
%   \item Distributed reference: \SI{62.5}{MHz} (HD and VD).
%   \item Timestamp alignment at the DTS firmware interface: \(\le\)\SI{5}{ns} within a module (requirement).
%   \item Endpoint recovered-clock quality: integrated jitter \(<\)\SI{180}{ps} (target).
% \end{itemize}
The timing endpoint physical layer shall follow the DUNE Timing System Endpoint
Integration Guide, \href{https://edms.cern.ch/document/2580777}{EDMS 2580777}.
The endpoint firmware interface is defined in
\href{https://edms.cern.ch/document/2395364}{EDMS 2395364 v2}. Timing endpoints
receive and transmit the DTS timing signal through commercial 1\,Gb/s SFP modules.
For final FD timing endpoints, the guide specifies 80\,km 1000BASE-BX bidirectional
SFPs; PDS DAPHNE endpoint SFPs should meet the following endpoint-side optical
requirements:
\begin{itemize}[nosep]
  \item PDS timing endpoints: TX \SI{1550}{nm}, RX \SI{1490}{nm}, as specified by EDMS 2580777.
  \item Associated timing-master optics: TX \SI{1490}{nm}, RX \SI{1550}{nm}.
  \item Connectors: LC on both ends; one LC fiber per endpoint (DAPHNE FEB and LCM).
  \item Fan-out: last stage on the DAQ mezzanine or mini-racks; PDS owns fibers from the last fan-out to endpoints.
  % \item Optical budget: minimum transmit power \(>\)-2\,dBm and minimum receive power \(<\)-24\,dBm.
  \item Endpoint SFP minimum transmit launch power \(\ge\)-2\,dBm and receiver sensitivity \(\le\)-24\,dBm; with \SI{17}{dB} maximum passive optical network loss, this gives at least \SI{5}{dB} worst-case margin.

  \item Digital Optical Monitoring (DOM) should be implemented.
  \item The SFP LOS signal shall be connected to the FPGA for standard endpoint firmware; the SFP I2C interface should be available for optical-status readout and debugging.
  \item The endpoint transmitter shall remain disabled except when requested by DTS; endpoint firmware shall control the SFP \texttt{TX\_DISABLE} input.
\end{itemize}




\paragraph{Startup, loss-of-lock, and recovery.}
\begin{itemize}[nosep]
  \item On power-up: endpoints wait for valid timing, report lock, then enable configuration. Spy-buffer and DAQ data transmission are enabled separately by the applicable control commands and run state.
  \item On loss-of-lock: flag packets as impacted; attempt auto-relock; report status via CCM; throttle/hold data if configured.
\end{itemize}


\subsection{Data Interface (10\,GbE Readout)}
\label{sec:physical-data}

\paragraph{Physical layer.}
\begin{itemize}[nosep]
  \item Optical: OM3/OM4 multimode; IEEE\,802.3ae 10GBASE-SR (\(\lambda\approx\)\SI{850}{nm}); LC connectors.
  \item Receiver input operating range: \(-17\)\,dBm to \(0\)\,dBm (nominal range used for link budgeting).
  \item One or more 10\,GbE lanes per DAPHNE FEB (design supports up to four 10\,GbE links if provisioned).
\end{itemize}

\paragraph{Link usage model.}
\begin{itemize}[nosep]
  \item Unidirectional data flow (DAPHNE\(\to\)DAQ). 
  \item DAQ provides the 10\,GbE/UDP transmitter firmware IP and the common DAQ \emph{header} block.
  \item Jumbo frames with MTU \(\sim\)9000\,B; UDP payloads aggregate multiple blocks to optimize efficiency.
\end{itemize}



\paragraph{Topology and scaling.}
\begin{itemize}[nosep]
  \item HD: \(\sim\)150 DAPHNE boards (40 ch per board).
  \item VD: \(\sim\)42 DAPHNE boards (32 ch per board).
\end{itemize}



\subsection{Control, Configuration, and Monitoring (CCM)}
\label{sec:ccm}

\paragraph{Physical/logical separation.}
\begin{itemize}[nosep]
  \item Each DAPHNE and each LCM provides a dedicated 1\,GbE SFP link reserved for CCM (no data readout on this link).
  \item DAQ and SC are authenticated OPC-UA clients of a shared PDS gateway. The gateway exposes a versioned, typed OPC-UA information model and is the sole routine client of the DAPHNE board services.
  \item The DAQ source configuration remains versioned in the DUNE-DAQ framework, including the DAPHNE schema in \texttt{daphnemodules}. The gateway implements a reviewed mapping from that model to OPC-UA variables and methods and then to the board's versioned Protobuf/ZMQ interface. Client applications shall not maintain divergent copies of the backend messages or connect directly to the board service.
  \item OPC-UA roles and method-level access control preserve the DAQ/SC ownership split. A DAQ method invocation is translated and executed by the neutral gateway; it is not forwarded to, approved by, or executed on behalf of DAQ by SC.
  \item Fibers up to the PDS mini-rack patch panels are provided by PDS; continuation and active switching are coordinated with DAQ and Technical Coordination.
\end{itemize}

\paragraph{Functional requirements.}
\begin{itemize}[nosep]
  \item DAQ owns the configuration that determines acquired data: timing endpoint and partition settings, AFE operating values, per-channel settings, trigger and spy-buffer selection, packetization, and readout links.
  \item SC owns equipment lifecycle and availability: board and AFE power enable, service lifecycle, controlled firmware deployment, management network, fans, rail and mezzanine power, environmental monitoring, and safe-state actions. PDS/DAQ change control owns qualification of firmware and its data-taking compatibility; SC deployment does not select an unapproved data-taking image.
  \item Health monitoring includes timing, data and CCM link state; temperature, voltage, current and fan status; service, firmware and bitstream identities; SFP DOM/LOS and I2C state; error counters; configuration readback; trigger state; and throttling statistics.
  \item The OPC-UA gateway and board service shall enforce disjoint write ownership, serialize accepted mutations, apply run-state guards, distinguish accepted, applied, and verified results, and give DPS inhibit or SC safe-state action priority over data-taking configuration.
\end{itemize}
